\section{Consequences and Levels of Evidence}
\label{sec:levels}

A feasible alternative is not yet a bug. In correct master--worker code, a wildcard with
two possible senders is the intended behaviour. \tool reports what the alternative
\emph{does}, and states what that claim rests on.

\subsection{Consequence classes}

\begin{table}[t]
\caption{Finding classes. ``Error'' classes count as bugs in the evaluation; the others are
warnings or information.}
\label{tab:classes}
\small
\begin{tabular}{@{}llp{0.58\linewidth}l@{}}
\toprule
class & level & meaning & counts as \\
\midrule
\cls{1a} & \Lone & the alternative message does not fit the receive (count or type signature) & error \\
\cls{1b} & \Lone & the alternative message belongs to another logical channel (tag) & warning \\
\cls{1f} & \Lone & same shape as the observed message; benign unless the continuation depends on it & info \\
\cls{1c} & \LoneSP, \Ltwo & a send is never received after the change & error \\
\cls{1d} & \Lone, \LoneSP, \Ltwo & deadlock after the change & error \\
\cls{1e} & \Ltwo & the continuation diverges from the trace, or the program aborts on its own check & error \\
\cls{2-buffering} & observed & the recorded run is realisable only above a threshold & error \\
\cls{2-collective} & observed & ranks disagree on a collective & error \\
\cls{2-blocked} & observed & the run ended with a rank blocked or spinning & error \\
\cls{probe-alt} & \Lone & a wildcard probe could have reported another message & warning \\
\cls{waitany-order} & \Lone & another request could have completed first & warning \\
\bottomrule
\end{tabular}
\end{table}

\Cref{tab:classes} lists the classes. Classes \cls{1a}, \cls{1b} and \cls{1f} follow from
the changed pair alone. A count or type mismatch is an error whatever the program does
next. The consequences that happen later, a message nobody receives or a deadlock,
depend on what the changed rank does after its receive. The horizon deliberately removes
exactly that part.

\subsection{Deadlocks inside the horizon}
\label{sec:l1dead}

Some deadlocks need no knowledge of the changed rank's continuation.

\begin{definition}[Contained deadlock]
Let $\Malt$ be feasible at $\Bt$ with horizon $\Hpre$. A set $D$ of ranks is a
\emph{contained deadlock} if no rank of $D$ is cut by \ref{H:changed}, every rank of $D$
is cut at an operation it has issued but cannot complete in $\Hpre$, and every send or
receive that could complete one of these operations belongs to a rank of $D$ and lies
after that rank's cut.
\end{definition}

\begin{theorem}[Deadlock at \Lone]
\label{thm:l1dead}
If $\Hpre$ contains a contained deadlock $D$, then for every
$\prog\in\Cons(\trace,\Hpre)$, every execution given by \cref{thm:sound} extends only to
executions in which the ranks of $D$ never complete their cut operations.
\end{theorem}

\noindent The ranks of $D$ are not cut by \ref{H:changed}, so their prefixes in $\Hpre$
are what every consistent program issues. Their blocked operations can only be completed
by operations of $D$, which never come. A wildcard receive belongs to a contained deadlock
only if every rank of its communicator is in $D$ or has terminated, because any other rank
might still send. This is the analogue of sync-preserving deadlock
prediction~\citep{tunc23spd} for message matching. \pending{Full proof.}

Every cut of the horizon descends from \ref{H:changed}, so a contained deadlock needs a
changed rank that stops \emph{before} its changed receive completes. Consider a nonblocking
wildcard \mpicall{Irecv(ANY)}, then \mpicall{Recv(from 1)}, then the wait. If the wildcard
takes rank~1's only message, the second receive has nothing left to take. The rank stops
there, before the wait whose result could change its behaviour. With a blocking wildcard,
the rank is cut right after it, and every rank that waits on it waits on an unknown
continuation. Such deadlocks remain \LoneSP or \Ltwo claims.

\subsection{Single-path consequences}

Most deadlocks and orphans an alternative causes do involve the changed rank. Earlier
predictors handle them by assuming that every rank issues the recorded calls whatever
it receives~\citep{forejt17mopper,huang20ase}. In our terms, they analyse $\SP(\trace)$.
At level \LoneSP, \tool explores $\SP(\trace)$ exhaustively with the changed match forced,
using a persistent-set reduction in which only wildcard matches and \mpicall{Waitany}
choices branch. It compares the outcomes with the unforced exploration, so that a hang
the recorded program has anyway is not attributed to the alternative. Each \LoneSP
finding states whether every execution that contains the changed pair has the
consequence or only some. Its pin fixes every wildcard receive of one witness execution.

\begin{lemma}[Single-path is unsound]
\label{lem:sp}
There is a trace $\trace$, an alternative $\Malt$, a threshold $\Bt$ and two programs
$\prog_1,\prog_2\in\Cons(\trace)$ such that every execution of $\prog_1$ at $\Bt$
realising $\Malt$ deadlocks, while $\prog_2$ has an execution at $\Bt$ realising $\Malt$
that terminates.
\end{lemma}

\noindent The master--worker program of \cref{sec:overview} at $\Bt=0$ is the witness.
Let $\prog_2$ be the real program, which replies to \texttt{status.MPI\_SOURCE}, and let
$\prog_1=\SP(\trace)$, which replies to the recorded destination. Both are consistent
with the recorded run. Under the alternative, $\SP(\trace)$ deadlocks and the real program
does not. A \LoneSP finding is therefore a claim about the recorded call sequence, not
about the program, and \tool labels it as such.

\subsection{Skeletons}

Level \Ltwo replaces $\SP(\trace)$ by a \emph{twin}: the recorded program with its
dependencies on message metadata made explicit. A destination can be computed from
\texttt{MPI\_SOURCE}. A branch can guard a collective. A payload can mirror the sender's
rank. The explorer over the twin decides \cls{1c}, \cls{1d} and \cls{1e} for programs
that agree with the twin. Twins are written by hand in this paper.
\pending{Static check that a continuation is metadata-independent, which makes \LoneSP
findings sound where it passes.}

\subsection{Confirmation by replay}
\label{sec:replay}

A theorem certifies the model; the tool can still diverge from it. Every receive and
probe finding therefore carries a pin. A pin names, per rank, a wildcard receive by its
ordinal and the sender it must take. The replay library rewrites that receive's source
before the call, translating world ranks into the receive's communicator. After the
receive completes, it logs the source the receive actually matched. A finding is
\emph{confirmed} when three conditions hold:
\begin{itemize}[leftmargin=*]
\item every pinned receive matched its named sender;
\item the run reaches the predicted outcome: a hang for \cls{1c} and \cls{1d}, the
  specific MPI error for \cls{1a};
\item a control replay that pins the observed matching finishes.
\end{itemize}
The first and the last condition matter because forcing an impossible match also
hangs. A pinned receive can only match its pinned source, so its completion, which the
library logs, is the evidence of the forced match.
